\section{Finite upper certificates and the boundary of affine branching}
\label{sec:xor-upper-affine}
The fixed-tolerance lower bounds have finite upper certificates of the
same exponential order in dimension. We give two, because they use
different information: the Bernstein certificate uses the order-two
preordering and polynomial graph identities, and the order-one
certificate for the quadratic formulation uses the exact quadratic hull
of the box.

\begin{proposition}[Order-two Bernstein certificates]
\label{prop:xor-bernstein-upper}
For every objective \eqref{eq:xor-objective} and $\epsilon>0$, at most
$\lceil3/\epsilon\rceil^n$ boxes certify the target $\OPT-\epsilon$
at order two. This holds both in the original variables and in the
quadratic formulation \eqref{eq:xor-factorable-formulation}; in the latter,
only original coordinates need subdivision. The quadratic moment hull
is not needed for this upper bound.
\end{proposition}
\begin{proof}
Set $M=\lceil3/\epsilon\rceil$ and partition each original interval
into $M$ equal intervals, of width $h=2/M\le2\epsilon/3$. Fix a resulting
box $B=\prod_i[a_i,a_i+h]$. Every clause cost has derivatives of
magnitude at most $1/2$ in each of its three coordinates. Moving between
points of its three-coordinate box one coordinate at a time changes its
value by at most $3h/2\le\epsilon$. Let $\beta_e$ be its minimum on
that box. In particular
\begin{equation}\label{eq:xor-grid-corner-bound}
 \frac1m\sum_e\beta_e\ge F(a)-3h/2\ge\OPT-\epsilon.
\end{equation}
For $i\in e$, write $\lambda_i^1=(x_i-a_i)/h$ and
$\lambda_i^0=(a_i+h-x_i)/h$. Multiaffine interpolation is the exact
polynomial identity
\begin{equation}\label{eq:xor-bernstein-identity}
 c_e(x)-\beta_e=
 \sum_{s\in\{0,1\}^e}\bigl(c_e(a_e+hs)-\beta_e\bigr)
       \prod_{i\in e}\lambda_i^{s_i}(x_i).
\end{equation}
Each coefficient is nonnegative, and each basis product is a positive
multiple of three coordinate-bound slacks. The minimum $\beta_e$ is
attained at a corner by multiaffinity. Thus the node preordering through
degree four certifies $L[c_e]\ge\beta_e$. Averaging and
\eqref{eq:xor-grid-corner-bound} give the original upper certificate.

For the lifted formulation, keep every auxiliary interval $[-1,1]$.
The exact identity
\begin{equation}\label{eq:bernstein-graph-transfer}
 v_e-x_ix_jx_k=(v_e-u_ex_k)+x_k(u_e-x_ix_j)
\end{equation}
has graph multipliers of total degree at most three. Order two therefore
implies $L[v_e]=L[x_ix_jx_k]$. The same Bernstein certificate in the
original bound slacks proves $L[\Phi]\ge\OPT-\epsilon$. There are
$M^n$ lifted boxes, and together they cover the full feasible graph.
\end{proof}

\begin{proposition}[Order-one certificates from the quadratic box hull]
\label{prop:xor-order-one-upper}
The quadratic formulation \eqref{eq:xor-factorable-formulation} admits
an order-one certificate at target $\OPT-\epsilon$ with at most
$\lceil3/\epsilon\rceil^n$ lifted boxes under the exact quadratic-box-hull
oracle, for every $\epsilon>0$. This certificate uses only the original
coordinate subdivisions and the two graph equations per clause imposed
on moments.
\end{proposition}
\begin{proof}
Use the same grid, with original lower corner $a$ and width $h=2/M$,
and leave auxiliaries in $[-1,1]$. Any feasible order-one functional
has its first and second moments realized by an actual distribution $\mu$
on this entire lifted box. The graph moment equations give
\begin{equation}\label{eq:order-one-graph-moments}
 L[u_e]=L[x_ix_j],\qquad L[v_e]=L[u_ex_k].
\end{equation}
We use these equalities only in expectation; $\mu$ need not satisfy the
graph equations pointwise. On the full box,
\[
 |x_ix_j-a_ia_j|
 \le |x_i-a_i|\,|x_j|+|a_i|\,|x_j-a_j|\le2h,
 \qquad |u_e(x_k-a_k)|\le h.
\]
Taking expectations, all of which are degree at most two, yields
\begin{align}
 |L[v_e]-a_ia_ja_k|
 &\le |L[u_ex_k]-a_kL[u_e]|
          +|a_k|\,|L[u_e]-a_ia_j|\nonumber\\
 &\le h+2h=3h.\label{eq:order-one-corner-estimate}
\end{align}
For either clause sign this gives
$L[(1-b_ev_e)/2]\ge c_e(a)-3h/2$.
Averaging proves $L[\Phi]\ge F(a)-3h/2\ge\OPT-\epsilon$.
The inequality holds for every feasible node functional and hence for
the infimum. The $M^n$ boxes cover the graph.

There is also an explicit quadratic inequality certifying each step. Set
\begin{equation}\label{eq:order-one-quadratic-cut}
 q_e=\frac{3h}{2}-\frac{b_e}{2}
       \bigl[u_e(x_k-a_k)+a_k(x_ix_j-a_ia_j)\bigr].
\end{equation}
The preceding pointwise bounds give $q_e\ge0$ on the \emph{full} lifted
box. Moreover,
\begin{equation}\label{eq:order-one-cut-identity}
 \frac{1-b_ev_e}{2}-c_e(a)+\frac{3h}{2}
 =q_e-\frac{b_e}{2}
       \bigl[(v_e-u_ex_k)+a_k(u_e-x_ix_j)\bigr].
\end{equation}
All polynomials in this identity have lifted degree at most two. The
quadratic hull gives $L[q_e]\ge0$, and the graph moment equations remove
the two remaining terms. The certificate thus uses a valid quadratic box
inequality linearly, without any higher-degree localizer or
graph-supported realizing distribution.
\end{proof}

At $\epsilon=1/16$, both propositions give $48^n$ regions. For the
family with $\OPT\ge1/8$, these also certify relative tolerance $1/2$
with the best incumbent: $\OPT-1/16\ge\OPT/2$. Consequently, for all sufficiently large $n$, the quadratic formulation
has an order-one lower bound $2^{7n/3072}$ and an order-one upper bound
$48^n$. These counts measure regions; we do not claim that the exact
quadratic-box-hull oracle is computationally tractable. For the original
cubic objective the order-two Bernstein certificate is the applicable
upper bound.

\subsection{Two obstructions to the same proof under affine branching}
The parity-rank count requires restricted coordinates to be signed
products with bounded support. A balanced affine halfspace behaves
differently, even under an actual probability distribution.

\begin{proposition}[A halfspace rejects the preserved quadratic moments]
\label{prop:affine-quadratic-obstruction}
For every $n\ge1$, let $\mathcal E$ be uniform expectation on
$\pmcube$ and put $S(x)=\sum_i x_i$. The first and second moments of
$\mathcal E$ have no actual realization on
\[
 H=\{x\in[-1,1]^n:S(x)\ge0\},
\]
although $H$ contains at least half the Boolean witnesses. Equivalently,
adding the affine coordinate $z=S/n$ and restricting it to $[0,1]$
rejects the preserved moments by the valid quadratic inequality $z^2\le z$.
\end{proposition}
\begin{proof}
The uniform moments give $\mathcal E[S]=0$ and
$\mathcal E[S^2]=n$. A distribution supported where $S\ge0$ with
mean $S=0$ must have $S=0$ almost surely, contradicting that second
moment. In the affine coordinate the same contradiction is
$\mathcal E[z]=0$ and $\mathcal E[z^2]=1/n>0$.
The involution $w\mapsto-w$ proves $\Prob\{S(w)\ge0\}\ge1/2$;
for even $n$ the boundary can add mass.
\end{proof}

\begin{proposition}[Substitution cost for a halfspace localizer]
\label{prop:affine-substitution-obstruction}
Start from the same uniform Boolean expectation, fix an arbitrary
coordinate set $C$ to arbitrary signs $w_C$, and leave the others
independently uniform. If the resulting functional $L$ satisfies
\[
 L[SP^2]\ge0\qquad(1+2\deg P\le2r)
\]
for an integer $r\ge1$, then necessarily
\begin{equation}\label{eq:affine-substitution-count}
 |C|\ge\min\{r-1,\lceil n/2\rceil\}.
\end{equation}
\end{proposition}
\begin{proof}
Write $t=|C|$, $a=\sum_{i\in C}w_i$, and $s=n-t$. Suppose both
$t<r-1$ and $t<\lceil n/2\rceil$. Then $t\le r-2$ and $s\ge t+1$.
If $a<0$, $P=1$ already violates the localizer. Otherwise $a$ is an
integer in $[0,t]$. Select $a+1$ unrestricted coordinates, which is
possible because $a+1\le t+1\le s$, and set
\[
 P(x)=\prod_{i\text{ selected}}\frac{1-x_i}{2}.
\]
Its degree is $a+1\le r-1$, so the localizer has permitted total degree.
On its indicator event the selected variables sum to $-(a+1)$, all
other unrestricted variables have conditional mean zero, and the fixed
ones sum to $a$. Boolean idempotence therefore gives the exact value
\[
 L[SP^2]=(a-(a+1))\,2^{-(a+1)}=-2^{-(a+1)}<0.
\]
This is the contradiction. For $r=1$ the conclusion is the stated trivial
bound $t\ge0$.
\end{proof}

At order proportional to $n$, this halfspace requires linearly many
substitutions in the displayed construction while retaining at least half
the witnesses. Thus, for general affine inequalities, a large substitution
cost does not imply exponentially small witness mass. These propositions
obstruct this proof method only. They do not rule out a different node
functional, nor show that the constant-gap XOR objective has a short
affine-branch certificate.

\subsection{Comparison with other disjunctive certificate models}
\label{subsec:xor-literature-boundary}
Disjunctive sum-of-squares certificates can keep polynomial degree fixed
while increasing the number of regions. Ahmadi, Dash, Hua, and Stellato
\citep[Sections 2--4 and 6]{AhmadiEtAl2026} develop such certificates
and a spatial framework on simplicial cones intersected with the sphere,
with a corresponding simplex method for matrix copositivity. Their
Section~6 states termination of the algorithms at positive prescribed
tolerance, which is compatible with exponentially many regions. The
coordinate-product bounds here do not apply to all simplicial or
algebraic disjunctions in that framework.

Stabbing Planes instead branches on an integer linear inequality
$ax\ge b$ and its integer negation $ax\le b-1$. It may discard the
intervening slab because the slab has no integer points
\citep[Section 1]{BeameEtAl2023}. Continuous subdivision must account
for those points and therefore has a different covering requirement.
Beame et al.\ give quasipolynomial-size Stabbing Planes refutations of
Tseitin formulas \citep[Theorem 1]{BeameEtAl2023}. Fleming et al.\ give
quasipolynomial-size Cutting Planes refutations for CNF encodings of
inconsistent finite-field linear systems, and size lower bounds for a
coefficient-restricted Stabbing Planes subsystem
\citep[Theorems 1.3--1.4]{FlemingEtAl2021}. These results neither transfer a parity SOS obstruction directly to
unrestricted affine branching nor give a constant-gap certificate for our
objective. A refutation of perfect Boolean satisfiability proves only
that some clause fails. Its direct consequence for the normalized
objective is $F\ge1/m$ on Boolean vertices, and hence on the cube by
multiaffinity; it does not give $F\ge1/16$ as $m$ grows. Comparable
constant-gap region bounds for a stronger affine-branching model remain
open here in both directions.
